\documentclass[11pt]{article}
\usepackage{mathunal-base}
\mathunalfuente{serif}

\renewcommand{\examtitulo}{Tercer Examen Parcial de Métodos Numéricos (20\%)}
\renewcommand{\examfecha}{Noviembre 20 de 2006}

\begin{document}

\examheader

\vspace{2pt}
\noindent Nombre: \dotfill\ Carné: \dotfill\ Grupo: \dotfill

\vspace{8pt}
\noindent\textbf{La interpretación del examen hace parte de la evaluación, por tanto no se admiten preguntas comenzada la prueba. Duración: 1 hora y 50 minutos.}

\vspace{10pt}
\noindent\textbf{Parte 1 Preguntas de selección múltiple. Cada una tiene un valor de 6\%.}

\noindent Escoja la única respuesta de cada pregunta y señálela haciendo una X sobre la letra correspondiente. Si en alguna pregunta señala más de una respuesta se le asigna un puntaje de cero a esa pregunta. Todas las preguntas valen lo mismo.

\begin{enumerate}
\item{} Sean $R=\{(t,y)\ /\ 0\leq t\leq 2,\ 0<y<3\}$ y $f(t,y)=e^{-2t}y$. De la lista siguiente el menor número que le sirve a $f(t,y)$ como constante de Lipschitz con respecto a su variable $y$ en $R$ es

\vspace{4pt}
\noindent a. $e^{-4}$ \hspace{1cm} b. $1$ \hspace{1cm} c. $3$ \hspace{1cm} d. $6e^{-4}$

\item{} Se desea resolver el PVI $y'=e^{-2t}y,\ y(0)=1$ utilizando el método de Taylor de orden 2. La fórmula de avance es:

\vspace{4pt}
\noindent a. $y_{k+1} = \left[1+he^{-2t_k}-\frac{h^2}{2}\left(e^{-4t_k}-2e^{-2t_k}\right)\right]y_k$

\vspace{4pt}
\noindent b. $y_{k+1} = \left[1+he^{-4t_k}-\frac{h^2}{2}\left(e^{-4t_k}-2e^{-2t_k}\right)\right]y_k$

\vspace{4pt}
\noindent c. $y_{k+1} = \left[1+he^{-2t_k}+\frac{h^2}{2}\left(e^{-4t_k}-2e^{-2t_k}\right)\right]y_k$

\vspace{4pt}
\noindent d. $y_{k+1} = \left[1+he^{-2t_k}-h^2e^{-2t_k}\right]y_k$

\item{} Utilizamos el método de Euler para aproximar el valor $y(0.4)$, donde $y(t)$ es la solución del PVI
\[
y' = t^2-y,\qquad y(0)=1. \tag{1}
\]
Si la aproximación que entrega el método es $\dfrac{81}{125}$, entonces se usó

\vspace{4pt}
\noindent a. $h=0.4$ \hspace{1cm} b. $h=0.2$ \hspace{1cm} c. $h=0.1$ \hspace{1cm} d. Ninguno de los anteriores

\item{} Para resolver el problema de contorno $y''+ty'-t^2y=2t^3,\ y(0)=1,\ y(1)=-1$ por el método del disparo, se deben resolver dos problemas de valor inicial, uno de los cuales es

\vspace{4pt}
\noindent a. $u''=-tu'+t^2u+2t^3,\ u(0)=1,\ u'(0)=-1$.

\noindent b. $v''=-tv'+t^2v,\ v(0)=0,\ v'(0)=1$.

\noindent c. $u''=-tu'+t^2u+2t^3,\ u(0)=0,\ u'(0)=1$.

\noindent d. $v''=-tv'+t^2v,\ v(0)=0,\ v'(0)=-1$.

\item{} Seguimos con el problema de contorno del ejercicio anterior. Supongamos que el PVI 1, el cual contiene el término $2t^3$, tiene solución $z_1(t)$ y que el PVI 2, que no contiene el término $2t^3$, tiene solución $z_2(t)$. Entonces la solución del problema de contorno es

\vspace{4pt}
\noindent a. $y(t) = z_1(t) - \dfrac{1+z_1(1)}{z_2(1)}z_2(t)$

\vspace{4pt}
\noindent b. $y(t) = z_1(t) + \dfrac{1-z_1(1)}{z_2(1)}z_2(t)$

\vspace{4pt}
\noindent c. $y(t) = z_2(t) - \dfrac{1+z_2(1)}{z_1(1)}z_2(t)$

\vspace{4pt}
\noindent d. $y(t) = z_2(t) - \dfrac{1+z_1(1)}{z_2(1)}z_1(t)$
\end{enumerate}

\vspace{10pt}
\noindent\textbf{Parte 2 Preguntas abiertas}

\begin{enumerate}
\setcounter{enumi}{5}
\item{} (20\%) Con respecto al PVI $y'=t^2-y,\ y(0)=1$, ¿Cuál es el valor que se obtiene con el método de Heun como aproximación para $y(0.4)$ utilizando $h=0.2$?

\vspace{4pt}
\noindent \textit{Nota: La fórmula de avance del método de Heun es $y_{k+1}=y_k+\frac{h}{2}\left[f(t_k,y_k)+f(t_k+h,y_k+hf(t_k,y_k))\right]$.}

\vspace{4cm}

\item{} (20\%) Se desea convertir el PVI de segundo orden $2y''-5y'-3y=45e^{2t}$, con $x(0)=2$ y $x'(0)=1$ a un PVI vectorial de la forma $Z'=F(t,Z)$ con $Z(0)=p$. Indique las expresiones para $Z$, $F(t,Z)$ y $p$.

\vspace{4cm}

\item{} Considere el problema
\[
\left\{
\begin{aligned}
u_t(x,t) &= u_{xx}(x,t) & 0<x<1,\ t>0\\
u(x,0) &= 1-\left|x-\tfrac{1}{2}\right| & 0\leq x\leq 1\\
u(0,t) &= 0 & t\geq 0\\
u(1,t) &= 0 & t\geq 0
\end{aligned}
\right.
\]
\begin{enumerate}
\item{} (10\%) ¿Será estable el método progresivo o explícito con los parámetros $h=0.1$ y $k=0.0125$? \textbf{Justifique.}

\vspace{4cm}

\item{} (20\%) Utilice el método progresivo o explícito para resolver el problema (2) con $h=0.25$ y $k=0.0125$ para completar los valores faltantes de la tabla denotadas $A$, $B$, $C$ y $D$.
\[
\begin{array}{c|ccccc}
 & x=0 & x=0.25 & x=0.5 & x=0.75 & x=1\\ \hline
t=0 & 0 & 0.75 & 1 & 0.75 & 0\\
t=0.0125 & 0 & 0.65 & A & B & 0\\
t=0.025 & 0 & C & 0.8 & 0.57 & D
\end{array}
\]

\vspace{4cm}
\end{enumerate}
\end{enumerate}

\examfooter
\end{document}
